% Macro x empresa e horizontes de previsão — experimento de 18/08/2026.
% Fonte: Mestrado_PETR4/macro_vs_empresa.json
% Script: src/modelagem/10_macro_vs_empresa_horizontes.py
\subsection{Que notícia informa o risco, e em que prazo?}
\label{sec:macro_empresa}

A Seção~\ref{sec:relevancia_volatilidade} concluiu que o índice de sentimento, embora
significativamente associado à volatilidade, não supera o modelo HAR na previsão fora da amostra.
Duas escolhas de desenho não foram, contudo, examinadas naquela ocasião, e ambas se revelaram
questionáveis à luz da literatura recente: o \emph{recorte temático} do corpus e o \emph{horizonte}
de previsão.

\subsubsection{A hipótese}

\citeonline{bodilsen_exploiting_2025}, no \textit{Journal of Applied Econometrics}, conduziram
exercício de desenho equivalente ao aqui reportado --- acréscimo do sentimento de notícias a modelos
de volatilidade realizada --- e chegaram a duas conclusões que esta dissertação não havia
investigado. A primeira é que o sentimento de notícia \textbf{específica da empresa} não acrescenta
informação ao que a própria volatilidade passada já captura. A segunda é que o sentimento de notícia
\textbf{macroeconômica} melhora as previsões de forma significativa, com ganho substancialmente
maior em \textbf{horizontes longos}.

A hipótese descreveria com precisão a situação desta pesquisa. O recorte adotado como padrão na
Seção~\ref{sec:relevancia_volatilidade} --- empresa somada ao mercado de petróleo --- é o mais
próximo da categoria ``específica da empresa'', e toda a avaliação foi conduzida a um único dia de
distância. Caso a hipótese se sustentasse aqui, ter-se-ia examinado a fatia menos informativa do
corpus, no prazo menos favorável. Cabia verificar.

\subsubsection{Desenho}

Cruzaram-se cinco recortes temáticos do corpus com três horizontes de previsão. Os recortes agrupam
a taxonomia descrita na Seção~\ref{sec:metodologia_ism}: \textsc{empresa} reúne as categorias de
Empresa e Governança; \textsc{petróleo} corresponde ao Mercado de Petróleo; \textsc{macro} reúne
Geopolítica, Sanções e Navegação, e Macro de Energia. Os horizontes são de um pregão, de cinco
pregões --- aproximadamente uma semana --- e de vinte e dois pregões, aproximadamente um mês, com o
alvo definido como a média do logaritmo da volatilidade no período.

O procedimento replica o da Seção~\ref{sec:relevancia_volatilidade}: volatilidade estimada por
\citeonline{parkinson_extreme_1980}, modelo HAR de \citeonline{corsi_simple_2009} como referência,
coeficientes reestimados a cada pregão sobre janela expansiva, e comparação por
\citeonline{diebold_comparing_1995} com erro-padrão robusto a autocorrelação. A
Tabela~\ref{tab:macro_empresa} apresenta o ganho percentual de erro quadrático sobre o HAR isolado.

\begin{table}[htpb]
\centering
\caption{Ganho percentual sobre o modelo HAR na previsão da volatilidade da PETR4, por recorte temático e horizonte (795 previsões fora da amostra). Valores positivos indicam que o sentimento reduz o erro.}
\label{tab:macro_empresa}
\begin{tabular}{l r c c c c}
\hline
\textbf{Recorte} & \textbf{Notícias} & \textbf{1 dia} & \textbf{5 dias} & \textbf{22 dias} & \textbf{Média} \\ \hline
\textbf{Empresa} & 71.003 & $+1{,}03\%$ & $+0{,}37\%$ & $\mathbf{+1{,}77\%}$ & $\mathbf{+1{,}06\%}$ \\
Empresa + petróleo & 120.792 & $+0{,}30\%$ & $-0{,}43\%$ & $+0{,}21\%$ & $+0{,}03\%$ \\
Petróleo & 55.910 & $-0{,}12\%$ & $-0{,}48\%$ & $-1{,}06\%$ & $-0{,}55\%$ \\
\textbf{Macro} & 76.438 & $-0{,}33\%$ & $\mathbf{-1{,}09\%}$ & $\mathbf{-1{,}79\%}$ & $\mathbf{-1{,}07\%}$ \\
Todas & 205.697 & $-0{,}30\%$ & $-1{,}93\%$ & $-2{,}45\%$ & $-1{,}56\%$ \\
\hline
\end{tabular}
\vspace{0.2cm}
{\small \\ Fonte: Elaborado pelo autor (\texttt{src/modelagem/10\_macro\_vs\_empresa\_horizontes.py}, 2026-08-18). Em negrito, os valores discutidos no texto.}
\end{table}

\subsubsection{Resultado: a hipótese inverte-se}

O padrão observado é o oposto do reportado por \citeonline{bodilsen_exploiting_2025}, e a inversão é
estatisticamente detectável.

O recorte macroeconômico \textbf{degrada} a previsão, e o faz de forma significativa nos dois
horizontes mais longos: $p = 0{,}0146$ em cinco pregões e $p = 0{,}0200$ em vinte e dois, com o
teste favorecendo o HAR isolado em ambos. Não se trata, portanto, de mera ausência de ganho, e sim
de prejuízo mensurável --- acrescentar sentimento macroeconômico ao modelo piora as previsões.

O recorte da empresa, inversamente, é o único que apresenta ganho médio positivo em todos os
horizontes, e alcança $+1{,}77\%$ em vinte e dois pregões, com $p = 0{,}0574$. É o resultado mais
próximo da superação do HAR obtido nesta pesquisa, ainda que aquém do limiar convencional de cinco
por cento.

Registre-se que \textbf{a metade da hipótese relativa ao horizonte confirma-se}: o melhor desempenho
do recorte da empresa manifesta-se no prazo de um mês, e não no de um dia. O que se inverte é
exclusivamente a proposição sobre o tipo de notícia.

\subsubsection{Interpretação}

Duas circunstâncias explicam a divergência, e ambas dizem respeito à diferença entre o objeto de
\citeonline{bodilsen_exploiting_2025} e o desta pesquisa.

A primeira é a composição do que se denomina ``macro''. Aqueles autores empregam notícia
macroeconômica \textbf{doméstica} aplicada a ações \textbf{norte-americanas}, e a macroeconomia dos
Estados Unidos é determinante direto do preço de um ativo negociado naquele mercado. O recorte macro
desta pesquisa é composto majoritariamente de \textbf{geopolítica internacional} --- 46.412 das
76.438 notícias ---, categoria que abrange o conflito na Ucrânia, sanções, decisões da OPEP e
tensões no Oriente Médio. Para um ativo brasileiro isolado, tal fluxo comporta-se predominantemente
como ruído, e a sua inclusão no modelo introduz variância sem informação correspondente.

A segunda diz respeito à natureza da PETR4. Trata-se de empresa de controle estatal, cujo risco
idiossincrático --- política de preços de combustíveis, intervenção do acionista controlador,
alterações na direção executiva, política de distribuição de dividendos --- responde por parcela da
variação substancialmente maior do que a observada em uma ação típica do índice S\&P~500. É
coerente, portanto, que a notícia da própria empresa carregue, neste caso, mais informação sobre o
risco futuro do que o noticiário macroeconômico.

\subsubsection{Uma tensão entre associação e previsão}

O resultado impõe qualificação ao da Seção~\ref{sec:relevancia_volatilidade}, e a tensão merece
registro explícito por constituir achado metodológico próprio.

Para a \textbf{associação} com a volatilidade, medida por correlação, o melhor recorte é o de
empresa somada ao mercado de petróleo, que supera o corpus integral em $23\%$ com $p = 0{,}001$.
Para a \textbf{previsão} fora da amostra, o melhor recorte é o de empresa e governança, ao passo que
o acréscimo do mercado de petróleo reduz o ganho médio de $+1{,}06\%$ para $+0{,}03\%$.

\textbf{Os dois critérios não premiam o mesmo corte do corpus.} A explicação provável reside na
distinção já estabelecida na Seção~\ref{sec:auditoria_noticia}: a correlação é sensível às caudas,
e a notícia do mercado de petróleo contribui justamente nos episódios extremos --- choques do
barril ---, ao passo que a previsão fora da amostra é dominada pelo comportamento do pregão típico,
em que aquele fluxo adiciona ruído. Trata-se de mais uma manifestação do efeito de cauda que
perpassa os resultados desta dissertação, e de advertência prática: \emph{a escolha do corpus deve
ser orientada pelo objetivo, e não presumida transferível entre tarefas}.
